\begin{afterword}
\summaryhead

The post continues the author's account of the young self's view of morality. In 2000, Eliezer2000 still held that building a superintelligence is right if anything is right, so that no one could have a justifiable conflict of interest over the Singularity. Then came a small thought: some people might want an AI not to kill them even if life is meaningless. The obvious dodge was that in that case nothing is right. But Eliezer2000 kept thinking: such people would have a motive to oppose the project now.

Three things kept the question alive: perfectionism, a wish not to wrong Brian Atkins, the Singularity Institute's funder, and a rule of honesty backed by the thought that deceivers get found out. So Eliezer2000 began to ask how to specify a fallback morality for an AI. This was the first technical thinking about morality, and it slowly made morality less mysterious. The lesson: actions screen off justifications, so what you do, not your reasons, decides what you learn. The author aims this at AI builders with reasons not to think about Friendly AI, and adds that a great truth may first appear as one small doubt.

\responsehead

The post is memoir with a moral, and the memoir matches the record where I could check it against the author's writings of the time. The 1999 ``Frequently Asked Questions about the Meaning of Life'' confesses ignorance of life's purpose, as the post says, and calls its argument ``a formal solution to the Meaning of Life,'' which the post recalls as an argument the young author ``fondly claimed to be `formal'.''

One step of the story was told differently at the time. In May 2001, on the SL4 mailing list, the author said that what ``pried me loose of pure objective morality'' was the possibility that an objective morality would have to be built rather than discovered, and \textsc{Creating Friendly AI} gives a similar first step. The post puts first a different thought: people's wish not to be killed, and the conflict of interest it creates. The two may be steps of one change, and both accounts arrive at planning for the case in which morality is arbitrary. But the account of what came first is not the one given in 2001.

The lesson, ``Actions screen off justifications,'' is also the moral of ``My Best and Worst Mistake'' a week earlier; this post applies it to a second case from the same history.

The post's last turn is to unnamed ``AI wannabes.'' No one is named. Those who say that the urgency of building AI leaves no time for the problem are described as people who ``just basically aren't interested in the problem.'' The test offered them, to look at actions ``stripped of all their justifications,'' comes with a ``maybe.'' The line about smart people ``skilled at defending beliefs they arrived at for unskilled reasons'' adapts one of Michael Shermer's.

In short: a memoir of the young author's first step toward Friendly AI that matches the author's 1999 writing where checked, though the author's 2001 account names a different first step.
\end{afterword}
